\exercisesection{\S~5}

\begin{enumerate}
\def\labelenumi{\arabic{enumi})}
\item
  The index of \(\operatorname{Aut}_{0}(\mathfrak{g})\) in \(\operatorname{Aut}(\mathfrak{g})\) is finite.
\item
  We have \(\mathrm{Aut}_e(\mathfrak{g}) = \mathrm{Aut}(\mathfrak{g})\) if \(\mathfrak{g}\) is splittable, simple, and of type \(\mathrm{G}_2, \mathrm{F}_4\) or \(\mathrm{E}_8\) .
\item
  Let \(\mathfrak{h}\) be a splitting Cartan subalgebra of \(\mathfrak{g}\) , \(\mathfrak{b}\) a Borel subalgebra of \((\mathfrak{g},\mathfrak{h})\) , \(\mathfrak{n} = [\mathfrak{b},\mathfrak{b}]\) and \(N = \exp \operatorname{ad}_{\mathfrak{g}}\mathfrak{n}\) . Then
\end{enumerate}

\[
\operatorname{Aut} (\mathfrak {g}) = \mathrm{N}. \operatorname{Aut} (\mathfrak {g}, \mathfrak {h}). \mathrm{N}.
\]

(Let \(s \in \operatorname{Aut}(\mathfrak{g})\) . Apply Prop. 10 of §3, no. 3 to \(\mathfrak{b} \cap s(\mathfrak{b})\) , then apply Chap. VII, §3, no. 4, Th. 3.)

\begin{enumerate}
\def\labelenumi{\arabic{enumi})}
\setcounter{enumi}{3}
\tightlist
\item
  Let \(\mathfrak{h}\) be a Cartan subalgebra of \(\mathfrak{g}\) , and \(s\) an element of \(\operatorname{Aut}(\mathfrak{g},\mathfrak{h})\) such that \(sH \neq H\) for all non-zero \(H\) in \(\mathfrak{h}\) . Show that \(s\) is of finite order. (Reduce to the case in which \(\mathfrak{h}\) is splitting, and choose an integer \(n \geq 1\) such that \(\varepsilon(s)^n = 1\) . Then there exists \(\varphi \in \mathrm{T}_{\mathrm{Q}}\) such that \(f(\varphi) = s^n\) . Let \(\sigma\) be the transpose of \(s|\mathfrak{h}\) . Show that \(1 + \sigma + \sigma^2 + \cdots + \sigma^{n-1} = 0\) , and deduce that \(s^{n^2} = 1\) .)
\end{enumerate}

¶ 5) a) Let \(a \in \operatorname{Aut}(\mathfrak{g})\) , and \(\mathfrak{n}\) the nilspace of \(a - 1\) . Show that the following conditions are equivalent:

\begin{enumerate}
\def\labelenumi{(\roman{enumi})}
\item
  \(\operatorname{Ker}(a - 1)\) is a Cartan subalgebra of \(\mathfrak{g}\) .
\item
  \(\mathfrak{n}\) is a Cartan subalgebra of \(\mathfrak{g}\) , and \(a\in \mathrm{Aut}_0(\mathfrak{g})\) .
\item
  \(\dim\mathfrak{n}=\operatorname{rk}(\mathfrak{g})\) and \(a\in\operatorname{Aut}_{0}(\mathfrak{g})\) .
\end{enumerate}

\begin{enumerate}
\def\labelenumi{\alph{enumi})}
\setcounter{enumi}{1}
\item
  Assume from now on that k is algebraically closed. Let V be a vector space, R a root system in V, \(T_{Q}\) the group \(\operatorname{Hom}(Q(R), k^{*})\) , n an integer \(\geq 1\) and \(T_{n}\) the subgroup of \(T_{Q}\) consisting of the elements whose order divides n.~Let \(\zeta\) be a primitive \(n\) th root of unity in \(k\) . For all \(H \in \mathrm{P}(\mathrm{R}^{\vee})\) , let \(\psi(H)\) be the element \(\gamma \mapsto \zeta^{\gamma(H)}\) of \(\mathrm{T}_{\mathrm{Q}}\) . The map \(\psi\) is a homomorphism from \(\mathrm{P}(\mathrm{R}^{\vee})\) to \(\mathrm{T}_n\) , with kernel \(n\mathrm{P}(\mathrm{R}^{\vee})\) . Let \(t \in \mathrm{T}_n\) , and let \(\mathrm{C}\) be an alcove in \(\mathrm{P}(\mathrm{R}^{\vee}) \otimes \mathbf{R}\) . There exists \(w \in \mathrm{W}(\mathrm{R})\) and \(H \in \mathrm{P}(\mathrm{R}^{\vee})\) such that \(\frac{1}{n} H \in \overline{\mathrm{C}}\) and \(\psi(wH) = t\) . (Use Chap. VI, §2, no. 1.)
\item
  Let \(\mathfrak{h}\) be a Cartan subalgebra of \(\mathfrak{g}\) , and \(\mathbb{R}\) the root system of \((\mathfrak{g},\mathfrak{h})\) . Let \(n,\zeta ,\psi\) be as in \(b\) ), and \(f\) as in no. 2. Let \(H\in \mathrm{P}(\mathrm{R}^{\vee})\) . The set of elements of \(\mathfrak{g}\) invariant under \(f(\psi (H))\) is \(\mathfrak{h}\oplus \sum_{\alpha \in \mathbb{R}'}\mathfrak{g}^\alpha\) , where \(\mathbb{R}'\) is the set of \(\alpha \in \mathbb{R}\) such that \(\alpha (H)\in n\mathbf{Z}\) , and \(f(\psi (H))\) satisfies the conditions of \(a\) ) if and only if \(\frac{1}{n} H\) belongs to an alcove.
\end{enumerate}

\(d\) ) Assume from now on that \(\mathfrak{g}\) is simple. Let \(\mathfrak{h}\) and \(\mathbb{R}\) be as in \(c\) ), \((\alpha_{1},\ldots ,\alpha_{l})\) a basis of \(\mathbb{R}\) , \(h\) the Coxeter number of \(\mathbb{R}\) , \(\zeta\) a primitive \(h\) th root of unity in \(k\) , and \(H\) the element of \(\mathfrak{h}\) such that \(\alpha_{i}(H) = 1\) for \(i = 1,\dots,l\) . Prove the following properties:

\begin{enumerate}
\def\labelenumi{(\roman{enumi})}
\item
  The homomorphism \(\gamma \mapsto \zeta^{\gamma (H)}\) from Q(R) to \(k^{*}\) defines an element of \(\operatorname {Aut}(\mathfrak{g},\mathfrak{h})\) which satisfies the conditions in \(a\) ), and has order \(h\) . (Use \(c\) ), Chap. VI, §2, Prop. 5 and Chap. VI, §1, Prop. 31.)
\item
  Every automorphism of \(\mathfrak{g}\) of finite order satisfies the conditions in \(a\) ) and is of order \(\geq h\) .
\item
  The automorphisms of \(\mathfrak{g}\) of order \(h\) satisfying the conditions of \(a\) ) form a conjugacy class in \(\mathrm{Aut}_e(\mathfrak{g})\) . (Use Prop. 5 of no. 3.)
\end{enumerate}

\begin{enumerate}
\def\labelenumi{\alph{enumi})}
\setcounter{enumi}{4}
\tightlist
\item
  Let \(\mathfrak{h}\) and \(\mathrm{R}\) be as in \(c\) ), and let \(w\) be a Coxeter transformation in \(\mathrm{W(R)}\) . Let \(s \in \operatorname{Aut}(\mathfrak{g}, \mathfrak{h})\) be such that \(\varepsilon(s) = w\) . Show that \(s\) satisfies the conditions in \(a\) ), and is of order \(h\) . (Use Chap. VI, §1, Prop. 33, Chap. V, §6, no. 2 and Chap. VII, §4, Prop. 9.)
\end{enumerate}

\(f)\) If \(s \in \operatorname{Aut}(\mathfrak{g})\) , the following conditions are equivalent:

\begin{enumerate}
\def\labelenumi{(\roman{enumi})}
\item
  \(s\) satisfies the conditions in \(a\) ), and is of order \(h\) ;
\item
  there exists a Cartan subalgebra \(\mathfrak{h}\) of \(\mathfrak{g}\) stable under \(s\) such that \(s|\mathfrak{h}\) is a Coxeter transformation in the Weyl group of \((\mathfrak{g},\mathfrak{h})\) . (Use \(d\) ) and \(e\) ).
\end{enumerate}

\(g\) ) The characteristic polynomial of the automorphism in \(d\) ) (i) is

\[
\mathrm{A} (\mathrm{T}) = (\mathrm{T} - 1) ^ {l} \prod_ {\alpha \in \mathrm{R}} (\mathrm{T} - \zeta^ {\alpha (H)}).
\]

That of the automorphism \(s\) in \(e\) ) is

\[
\mathrm{B} (\mathrm{T}) = \left(\mathrm{T} ^ {h} - 1\right) ^ {l} \prod_ {i = 1} ^ {l} \left(\mathrm{T} - \zeta^ {m _ {i}}\right) \quad \left(m _ {i} \text { being   the   exponents   of } \mathrm{R}\right).
\]

(Use Prop. 33 (iv) of Chap. VI, §1, no. 11.) Deduce from the relation \(\mathrm{A(T)} = \mathrm{B(T)}\) that, for all \(j \geq 1\) , the number of \(i\) such that \(m_i \geq j\) is equal to the number of \(\alpha\in R_{+}\) such that \(\alpha(H)=j\) ; hence recover the result of Exerc. 6 c) of Chap. VI, §4. \(^{12}\)

\begin{enumerate}
\def\labelenumi{\arabic{enumi})}
\setcounter{enumi}{5}
\item
  Assume that \(k = \mathbf{R}\) or \(\mathbf{C}\) . Let \(G\) be the Lie group \(\mathrm{Aut}_0(\mathfrak{g})\) . Show that an element \(a \in G\) is regular (in the sense of Chap. VII, §4, no. 2) if and only if it satisfies the conditions of Exerc. 5 a).
\item
  Assume that g is splittable. Let \(\mathrm{B}(\mathfrak{g})\) be the canonical basis of the canonical Cartan subalgebra of g (no. 3, Remark 2). If \(s \in \mathrm{Aut}(\mathfrak{g})\) , s induces a permutation of \(\mathrm{B}(\mathfrak{g})\) ; denote by \(\operatorname{sgn}(s)\) the sign of this permutation. Show that s operates on \(\bigwedge^{n} g\) (with \(n = \dim g\) ) by
\end{enumerate}

\[
x \mapsto \operatorname{sgn} (s). x.
\]

¶8) Let \(\bar{k}\) be an algebraic closure of k, and \(\bar{g} = \bar{k} \otimes_{k} g\) . The Galois group \(\operatorname{Gal}(\bar{k}/k)\) operates naturally on \(\bar{g}\) , on the canonical Cartan subalgebra of \(\bar{g}\) (no. 3, Remark 2), as well as on its root system \(\bar{R}\) and its canonical basis \(\bar{B}\) . We thus obtain a continuous homomorphism (i.e.~with open kernel)

\[
\pi : \operatorname{Gal} (\bar {k} / k) \to \operatorname{Aut} (\bar {\mathrm{R}}, \bar {\mathrm{B}}).
\]

Show that \(\mathfrak{g}\) is splittable if and only if the following two conditions are satisfied:

\begin{enumerate}
\def\labelenumi{(\roman{enumi})}
\item
  The homomorphism \(\pi\) is trivial.
\item
  g has a Borel subalgebra b.
\end{enumerate}

(Show that a Cartan subalgebra \(\mathfrak{h}\) contained in \(\mathfrak{b}\) is splitting if and only if \(\pi\) is trivial.)

\begin{enumerate}
\def\labelenumi{\arabic{enumi})}
\setcounter{enumi}{8}
\tightlist
\item
  Let R be a reduced root system, B a basis of R, and \((\mathfrak{g}_{0}, \mathfrak{h}_{0}, \mathrm{B}, (X_{\alpha})_{\alpha \in \mathrm{B}})\) a corresponding framed semi-simple Lie algebra (§4). Let \(\bar{k}\) be an algebraic closure of k, and \(\rho : \operatorname{Gal}(\bar{k}/k) \to \operatorname{Aut}(\mathrm{R}, \mathrm{B})\) a continuous homomorphism (cf.~Exerc. 8); if \(\sigma \in \operatorname{Gal}(\bar{k}/k)\) , denote by \(\rho_{\sigma}\) the k-linear automorphism of \(\bar{g}_{0} = \bar{k} \otimes_{k} g_{0}\) such that \(\rho_{\sigma}(X_{\alpha}) = X_{\rho(\sigma)\alpha}\) . On the other hand, the natural operation of \(\operatorname{Gal}(\bar{k}/k)\) on \(\bar{k}\) can be extended to an operation on \(\bar{g}_{0}\) . Let g be the subset of \(\bar{g}_{0}\) consisting of the elements x such that \(\rho_{\sigma}(x) = \sigma^{-1}.x\) for all \(\sigma \in \operatorname{Gal}(\bar{k}/k)\) .
\end{enumerate}

\begin{enumerate}
\def\labelenumi{\alph{enumi})}
\item
  Show that \(\mathfrak{g}\) is a \(k\) -Lie subalgebra of \(\bar{\mathfrak{g}}_0\) , and that the injection of \(\mathfrak{g}\) into \(\bar{\mathfrak{g}}_0\) extends to an isomorphism from \(\bar{k} \otimes_k \mathfrak{g}\) to \(\bar{\mathfrak{g}}_0\) . In particular, \(\mathfrak{g}\) is semi-simple.
\item
  Let \(\mathfrak{b}_0\) be the subalgebra of \(\mathfrak{g}_0\) generated by \(\mathfrak{h}_0\) and the \(X_{\alpha}\) . Put
\end{enumerate}

\[
\bar {\mathfrak {b}} _ {0} = \bar {k} \otimes_ {k} \mathfrak {b} _ {0}, \quad \mathfrak {h} = \mathfrak {g} \cap \bar {\mathfrak {b}} _ {0}, \quad \bar {\mathfrak {h}} _ {0} = \bar {k} \otimes_ {k} \mathfrak {h} _ {0}, \quad \mathfrak {h} = \mathfrak {g} \cap \bar {\mathfrak {h}} _ {0}.
\]

Show that \(\bar{k}\otimes_{k}b=\bar{b}_{0}\) and that \(\bar{k}\otimes_{k}h=\bar{h}_{0}\) , so that b is a Borel subalgebra of g and h is a Cartan subalgebra contained in it.

\begin{enumerate}
\def\labelenumi{\alph{enumi})}
\setcounter{enumi}{2}
\tightlist
\item
  Show that the homomorphism \(\pi\) associated to the algebra g (cf.~Exerc. 8) is equal to \(\rho\) .
\end{enumerate}

¶ 10) Let h be a splitting Cartan subalgebra of g, and \((X_{\alpha})_{\alpha\in\mathbb{R}}\) a Chevalley system in \((\mathfrak{g},\mathfrak{h})\) , cf.~§2, no. 4. If \(\alpha\in R\) , put

\[
\theta_ {\alpha} = e ^ {\mathrm{ad} X _ {\alpha}} e ^ {\mathrm{ad} X _ {- \alpha}} e ^ {\mathrm{ad} X _ {\alpha}},
\]

and denote by \(\overline{W}\) the subgroup of \(\operatorname{Aut}_{e}(\mathfrak{g},\mathfrak{h})\) generated by the \(\theta_{\alpha}\) .

\begin{enumerate}
\def\labelenumi{\alph{enumi})}
\item
  Show that \(\varepsilon (\overline{\mathrm{W}}) = \mathrm{W}(\mathrm{R})\) .
\item
  Let \(s \in \overline{\mathrm{W}}\) , and let \(w = \varepsilon(s)\) . Show that
\end{enumerate}

\[
s (X _ {\alpha}) = \pm X _ {w (\alpha)} \quad \text {   for   all   } \alpha \in \mathrm{R}.
\]

(Use Exerc. 5 of §2.)

\begin{enumerate}
\def\labelenumi{\alph{enumi})}
\setcounter{enumi}{2}
\tightlist
\item
  Let M be the kernel of \(\varepsilon : \overline{\mathrm{W}} \to \mathrm{W}(\mathrm{R})\) . Show that M is contained in the subgroup of \(f(\mathrm{T}_{\mathrm{Q}})\) consisting of the elements \(f(\varphi)\) such that \(\varphi^2 = 1\) . Show that M contains the elements \(f(\varphi_{\alpha})\) defined by \(\varphi_{\alpha}(\beta) = (-1)^{\langle \beta, \alpha^{\vee} \rangle}\) (remark that \(\theta_{\alpha}^{2} = f(\varphi_{\alpha})\) ).
\end{enumerate}

\(d)\) *Let \(\varphi\in\mathrm{Hom}(\mathbf{Q},\{\pm1\})\) . Show that \(f(\varphi)\) belongs to M if and only if \(\varphi\) extends to a homomorphism from P to \(\{\pm1\}\) . (Sufficiency follows from the fact that M contains the \(f(\varphi_{\alpha})\) . To prove necessity, reduce to the case in which \(k=\mathbf{Q}\) , and use the fact that M is contained in \(f(\mathrm{T}_{\mathrm{Q}})\cap\mathrm{Aut}_{e}(\mathfrak{g})=\mathrm{Im}(\mathrm{T}_{\mathrm{P}})\) , cf.~§7, Exerc. 26 \(d)\) .) Deduce that M is isomorphic to the dual of the group \(\mathrm{Q}/(\mathrm{Q}\cap2\mathrm{P}).^{13}\)

\begin{enumerate}
\def\labelenumi{\arabic{enumi})}
\setcounter{enumi}{10}
\tightlist
\item
  With the notations of no. 2, assume that \(k\) is non-discrete and locally compact, hence isomorphic to \(\mathbf{R}\) , \(\mathbf{C}\) or a finite extension of \(\mathbf{Q}_p\) (Commutative Algebra, Chap. VI, §9, no. 3). For all \(n \geq 1\) , the quotient \(k^* / k^{*n}\) is finite (cf.~Commutative Algebra, Chap. VI, §9, Exerc. 3 for the ultrametric case). Deduce that the quotients \(T_Q / \text{Im}(T_P)\) and \(\text{Aut}(\mathfrak{g}) / \text{Aut}_e(\mathfrak{g})\) are finite.
\end{enumerate}

When \(k = \mathbf{R}\) , show that \(\mathrm{T_Q / Im(T_P)}\) is isomorphic to the dual of the \(\mathbf{F}_{2}\) -vector space \((\mathrm{Q} \cap 2\mathrm{P}) / 2\mathrm{Q}\) . When \(k = \mathbf{C}\) , we have \(\mathrm{T_Q = Im(T_P)}\) ; this is the integral subgroup of the Lie group \(\operatorname{Aut}(\mathfrak{g})\) with Lie algebra \(\mathfrak{h}\) .

¶ 12) Let \(\mathfrak{h}\) be a splitting Cartan subalgebra of \(\mathfrak{g}\) , A a subset of \(\mathfrak{h}\) , and \(s \in \operatorname{Aut}(\mathfrak{g})\) such that \(s\mathrm{A} = \mathrm{A}\) . Show that there exists \(t \in \operatorname{Aut}(\mathfrak{g}, \mathfrak{h})\) such that \(t|\mathrm{A} = s|\mathrm{A}\) and \(ts^{-1} \in \operatorname{Aut}_e(\mathfrak{g})\) . (Let \(\mathfrak{a}\) be the commutant of A in \(\mathfrak{g}\) ; this is a reductive subalgebra of \(\mathfrak{g}\) , of which \(sh\) and \(\mathfrak{h}\) are splitting Cartan subalgebras; deduce the existence of \(u \in \operatorname{Aut}_e(\mathfrak{a})\) such that \(ush = \mathfrak{h}\) ; show that there exists \(v \in \operatorname{Aut}_e(\mathfrak{g})\) extending \(u\) such that \(v|\mathrm{A} = \mathrm{Id}_{\mathrm{A}}\) ; take \(t = vs.\) ) Deduce that, if \(s \in \operatorname{Aut}_0(\mathfrak{g})\) , there exists \(w \in \mathrm{W}(\mathrm{R})\) such that \(w|\mathrm{A} = s|\mathrm{A}\) .

¶13) Let \((\mathfrak{g},\mathfrak{h},\mathrm{B},(X_{\alpha})_{\alpha\in\mathrm{B}})\) be a framed semi-simple Lie algebra, R the root system of \((\mathfrak{g},\mathfrak{h})\) , \(\Delta\) the corresponding Dynkin graph, and \(\varPhi\) a subgroup of

\[
\operatorname{Aut} (\mathrm{R}, \mathrm{B}) = \operatorname{Aut} (\Delta).
\]

If \(s \in \varPhi\) , extend \(s\) to an automorphism of \(\mathfrak{g}\) by the conditions

\[
s (X _ {\alpha}) = X _ {s \alpha}, \quad \text { and } \quad s (H _ {\alpha}) = H _ {s \alpha} \quad \text { for   all } \alpha \in \mathrm{B}, \text { cf.   Prop. } 1.
\]

We thus identify \(\varPhi\) with a subgroup of \(\operatorname{Aut}(\mathfrak{g},\mathfrak{h})\) ; denote by \(\tilde{\mathfrak{g}}\) (resp. \(\tilde{\mathfrak{h}}\) ) the subalgebra of g (resp. h) consisting of the elements invariant under \(\varPhi\) .

\begin{enumerate}
\def\labelenumi{\alph{enumi})}
\tightlist
\item
  Let \(\alpha \in \mathbf{B}\) , and let \(\mathbf{X} = \varPhi.\alpha\) . Show, by using the Plates in Chap. VI, that only two cases are possible:
\end{enumerate}

\begin{enumerate}
\def\labelenumi{(\roman{enumi})}
\item
  every element of \(X\) distinct from \(\alpha\) is orthogonal to \(\alpha\) ;
\item
  there exists a unique element \(\alpha'\) of \(X - \{\alpha\}\) that is not orthogonal to \(\alpha\) , and \(n(\alpha, \alpha') = n(\alpha', \alpha) = -1\) .
\end{enumerate}

\begin{enumerate}
\def\labelenumi{\alph{enumi})}
\setcounter{enumi}{1}
\tightlist
\item
  Let i be the restriction map \(h^{*} \to \tilde{h}^{*}\) , and let \(\tilde{\mathbf{B}} = i(\mathbf{B})\) ; the map \(\mathbf{B} \to \tilde{\mathbf{B}}\) identifies \(\tilde{B}\) with \(B/\varPhi\) . Show that \(\tilde{g}\) is a semi-simple Lie algebra, that \(\tilde{h}\) is a splitting Cartan subalgebra of it, and that \(\tilde{B}\) is a basis of \(\mathrm{R}(\tilde{\mathfrak{g}}, \tilde{\mathfrak{h}})\) . (Observe that \(\tilde{B}\) is contained in \(\mathrm{R}(\tilde{\mathfrak{g}}, \tilde{\mathfrak{h}})\) , and that every element of \(\mathrm{R}(\tilde{\mathfrak{g}}, \tilde{\mathfrak{h}})\) is a linear combination of elements of \(\tilde{B}\) with integer coefficients of the same sign.) If \(\tilde{\alpha} \in \tilde{B}\) , the corresponding inverse root \(H_{\tilde{\alpha}} \in \tilde{h}\) is given by
\end{enumerate}

\[
\begin{array}{l} H _ {\tilde {\alpha}} = \sum_ {i (\alpha) = \tilde {\alpha}} H _ {\alpha} \quad \text {in case (i) of a)} \\ H _ {\tilde {\alpha}} = 2 \sum_ {i (\alpha) = \tilde {\alpha}} H _ {\alpha} \quad \text {in case (ii) of a)}, \end{array}
\]

where the summation is over those elements \(\alpha \in \mathrm{B}\) such that \(i(\alpha) = \tilde{\alpha}\) . If \(\beta \in \mathrm{B}\) has image \(\tilde{\beta} \in \tilde{\mathrm{B}}\) , then

\[
  n (\tilde {\beta}, \tilde {\alpha}) = \sum_ {i (\alpha) = \tilde {\alpha}} n (\beta , \alpha) \quad \text { in   case   (i) }
\]

\[
n (\tilde {\beta}, \tilde {\alpha}) = 2 \sum_ {i (\alpha) = \tilde {\alpha}} n (\beta , \alpha) \quad \text { in   case   (ii) }.
\]

Deduce how the Dynkin graph of \(\mathrm{R}(\tilde{\mathfrak{g}},\tilde{\mathfrak{h}})\) is determined starting from the pair \((\varDelta,\varPhi)\) .

\begin{enumerate}
\def\labelenumi{\alph{enumi})}
\setcounter{enumi}{2}
\tightlist
\item
  Show that, if \(\mathfrak{g}\) is simple, so is \(\tilde{\mathfrak{g}}\) .
\end{enumerate}

If \(\mathfrak{g}\) is of type \(A_l\) , \(l \geq 2\) , and \(\varPhi\) is of order 2, then \(\tilde{\mathfrak{g}}\) is of type \(B_{l/2}\) if \(l\) is even, and of type \(C_{(l+1)/2}\) if \(l\) is odd.

If \(\mathfrak{g}\) is of type \(D_l\) , \(l \geq 4\) , and \(\varPhi\) is of order 2, then \(\tilde{\mathfrak{g}}\) is of type \(B_{l-1}\) .

If \(\mathfrak{g}\) is of type \(\mathrm{D}_4\) , and \(\varPhi\) is of order 3 or 6, then \(\tilde{\mathfrak{g}}\) is of type \(\mathrm{G}_2\) .

If g is of type \(E_{6}\) , and \(\Phi\) is of order 2, then \(\tilde{g}\) is of type \(F_{4}\) .

